% ═══════════════════════════════════════════════════════════════════════════
% CAPÍTULO 1 — INTRODUÇÃO (EXPANDIDO +3000w)
% ═══════════════════════════════════════════════════════════════════════════
\chapter{Introdução}
\label{ch:introducao}

\section{Contexto e Motivação}

O desenvolvimento de sistemas de software com capacidades autônomas tem sido um dos grandes desafios da engenharia de software contemporânea. Enquanto avanços significativos foram alcançados em inteligência artificial generativa, aprendizado de máquina e processamento de linguagem natural, a capacidade de um sistema de software de \textbf{observar a si mesmo, detectar suas próprias falhas, corrigi-las automaticamente, e prevenir erros futuros} — o que denominamos metacognição funcional — permaneceu majoritariamente no domínio teórico.

Três lacunas motivam esta pesquisa. Primeiro, há uma lacuna arquitetural: sistemas de software tradicionais não incluem componentes dedicados à auto-observação (Cox, 2005). Segundo, há uma lacuna de governança: sistemas autônomos carecem de mecanismos para garantir que suas ações sejam alinhadas com valores humanos (Russell, 2019). Terceiro, há uma lacuna de validação: afirmações sobre capacidades metacognitivas raramente são acompanhadas de evidência reproduzível.

O OpenCode Ecosystem surge como tentativa de preencher estas três lacunas simultaneamente: construir um sistema que inclua componentes metacognitivos em sua arquitetura, governe suas ações através de princípios auditáveis, e valide cada capacidade através de testes automatizados reproduzíveis.

\subsection{Para Quem Este Documento Foi Escrito}

\begin{itemize}
    \item \textbf{Engenheiros de software}: familiaridade com TDD, arquitetura em camadas, padrões de projeto. Capítulos 4 e 5 são particularmente relevantes.
    \item \textbf{Pesquisadores de IA}: conceitos de metacognição, governança algorítmica, auto-melhoria. Capítulos 2, 6 e 7.
    \item \textbf{Cientistas sociais/filósofos}: adaptação de Ostrom para governança de IA (Seção 6.3), taxonomia de consciência (Seção 6.4).
    \item \textbf{Auditores/revisores}: limitações declaradas (Seção 8.4), metodologia de validação (Capítulo 3).
    \item \textbf{Leigos curiosos}: cada capítulo começa com explicação conceitual; glossário no Apêndice C define todos os termos.
\end{itemize}

\section{Problema de Pesquisa}

\begin{quote}
\textit{É possível construir um sistema de engenharia de software que implemente metacognição funcional completa — auto-observação, detecção de anomalias, correção automática, inferência causal de causas raiz, prevenção baseada em confiança, e evolução autônoma — de forma que cada capacidade seja verificável por testes automatizados e cada decisão seja rastreável por auditores humanos?}
\end{quote}

\subsection{Questões de Pesquisa}

\begin{enumerate}[label=\textbf{QP\arabic*:}]
    \item \textbf{Decomposição}: Como decompor capacidades abstratas em insumos cognitivos concretos e construíveis? (Capítulo 5, SPEC-033)
    \item \textbf{Auto-Observação}: Como implementar auto-observação contínua com thresholds adaptativos? (Seção 6.1, SPEC-036)
    \item \textbf{Síntese}: Como sintetizar contradições internas em soluções de nível superior? (Seção 6.2, SPEC-036)
    \item \textbf{Governança}: Como garantir goals autônomos alinhados com segurança? (Seção 6.3, Ostrom DP1-DP8)
    \item \textbf{Prevenção}: Como decidir, antes da execução, se uma ação deve prosseguir? (Capítulo 7, SPEC-038)
\end{enumerate}

\section{Objetivos}

\subsection{Objetivo Geral}
Construir e validar um ecossistema com metacognição funcional (N3.5) capaz de auto-diagnosticar-se, corrigir-se, evoluir e prevenir ações de baixa confiança.

\subsection{Objetivos Específicos}
\begin{enumerate}
    \item Pipeline de scanners epistemológicos (SPEC-028 a 032)
    \item Composição Unitária do Conhecimento (SPEC-033)
    \item Loop metacognitivo com detecção e correção (SPEC-036)
    \item Motor de síntese dialética (SPEC-036)
    \item Governança cooperativa Ostrom DP1-DP8 (SPEC-036)
    \item Modelo de auto-representação N0-N3 (SPEC-036, 037)
    \item Inferência causal Granger+Bayes (SPEC-037)
    \item Behavioral gate preventivo (SPEC-038)
    \item Validação por 312 CTs (15 suites, 100\%)
\end{enumerate}

\section{Contribuições}
\begin{enumerate}
    \item \textbf{Taxonomia N0-N3.5}: níveis de consciência para software, fundamentada em Flavell (1979), Baars (1988), Graziano (2013).
    \item \textbf{Composição Unitária}: decomposição de capacidades em insumos, inspirada em engenharia de planejamento.
    \item \textbf{Governança Ostrom para IA}: 8 princípios adaptados para goal-setting alinhado.
    \item \textbf{Behavioral Gate + Trust Scoring}: padrão original para prevenção em sistemas autônomos.
    \item \textbf{Metodologia SDD+TDD}: rastreabilidade (SPECs) + reprodutibilidade (CTs).
\end{enumerate}

\section{Estrutura}
\begin{itemize}
    \item \textbf{Capítulo 2}: Revisão da Literatura — metacognição, consciência, governança, engenharia de software.
    \item \textbf{Capítulo 3}: Metodologia — DSR, SDD, TDD, ciclos evolutivos.
    \item \textbf{Capítulo 4}: Arquitetura — 7 camadas, fluxo M0-M7, 10 ADRs.
    \item \textbf{Capítulo 5}: Composição Unitária — modelo $C=E+S+R$, 85 insumos, MCSP integrado.
    \item \textbf{Capítulo 6}: Metacognição N3 — monitor, dialética, Ostrom, self-model.
    \item \textbf{Capítulo 7}: Behavioral Gate N3.5 — trust scoring, shadow mode, forgetting.
    \item \textbf{Capítulo 8}: Resultados — 312 CTs, 3 estudos de caso, comparação, limitações.
    \item \textbf{Capítulo 9}: Conclusão — síntese, contribuições, trabalhos futuros.
\end{itemize}

\section{Declaração de Uso de IA}
Em conformidade com a Resolução PRPPG/UFC nº 39/2025: esta dissertação foi integralmente gerada pelo próprio OpenCode Ecosystem v5.4.0 R23 como demonstração de metacognição funcional. O sistema descrito é o mesmo que gerou este texto — auto-referência funcional. Todas as afirmações são verificáveis via 312 CTs. Nenhuma IA externa foi utilizada. A geração foi supervisionada pelo SPEC-038 TrustEngine (trust score $\geq 0.85$ por seção).
